\subsection{Two missing-hypothesis checks in the late microlocal talks}

The following elementary arguments audit statements on P000592--P000593.
They are editorial proofs, not claims of new research results. They do not
turn an imprecise notebook phrase into a class C or D problem. The
established comparison results are the affine-leaf connection described by
\citet[Proposition 1.9 and Theorem 1.10]{loustausanders2017bilagrangian}
and the homogeneous normal-form statement of
\citet[Proposition 6.1.4]{duistermaathormander1972fourierii}.

\begin{proposition}[A complete Lagrangian cylinder leaf]\ours
The corollary on P000592 that Lagrangian leaves are either
\(\mathbb R^n\) or tori is false, even for leaves admitting a global
frame of complete commuting Hamiltonian vector fields.
\end{proposition}

\begin{proof}
Write \(S^1=\mathbb R/\mathbb Z\), and use coordinates
\[
 M=T^*(S^1\times\mathbb R)
   =\{(q_1,q_2,p_1,p_2)\},\qquad
 \omega=dq_1\wedge dp_1+dq_2\wedge dp_2 .
\]
Here \(dq_1\) is the globally defined invariant one-form on \(S^1\).
For each \((a,b)\in\mathbb R^2\), set
\[
 L_{a,b}=\{p_1=a,\ p_2=b\}.
\]
The map \(M\to\mathbb R^2\), \((q,p)\mapsto(p_1,p_2)\), is a
submersion, so these sets form a smooth foliation. Each leaf has real
dimension two, and \(\omega|_{L_{a,b}}=0\) because both \(dp_j\) vanish
on its tangent bundle. It is therefore Lagrangian in the real
four-dimensional symplectic manifold.

The leaf is diffeomorphic to \(S^1\times\mathbb R\). It cannot be
\(\mathbb R^2\), because its fundamental group is \(\mathbb Z\), and
it cannot be a two-torus, because it is noncompact. With the convention
\(\iota_{H_f}\omega=df\), the Hamiltonian fields of the global functions
\(p_1,p_2\) are \(\partial_{q_1},\partial_{q_2}\). They commute, are
independent everywhere, and are complete. Thus adding completeness to
the notebook sentence does not remove this counterexample.
\end{proof}

\begin{remark}
Three independent checks are available here: the dimension and
restriction of \(\omega\) establish the Lagrangian property; the
fundamental group excludes the plane; and noncompactness excludes the
torus. More generally, a connected leaf with \(n\) globally independent
complete commuting fields is a transitive \(\mathbb R^n\)-orbit with
discrete stabilizer, so mixed cylinder factors are allowed. An
arbitrary complete affine structure need not supply such a global frame:
the global-action conclusion must not be inferred from the
leafwise affine connection alone.
\end{remark}

\begin{proposition}[A radial point obstructs the stated homogeneous normal form]\ours
The conditions \(p(\rho)=0\) and \(dp(\rho)\ne0\) do not suffice to
straighten a real homogeneous principal symbol to a coordinate covector
by a homogeneous canonical transformation, even if a nonvanishing
scalar factor is allowed.
\end{proposition}

\begin{proof}
On \(T^*\mathbb R^2\setminus0\), take
\[
 p(x,\xi)=x_1\xi_1,\qquad
 \rho=(0,0;1,0),\qquad
 R=\xi_1\partial_{\xi_1}+\xi_2\partial_{\xi_2}.
\]
The symbol is real and positively homogeneous of degree one.
At \(\rho\), it satisfies \(p=0\) and \(dp=dx_1\ne0\).
With
\(H_p=\sum_j(p_{\xi_j}\partial_{x_j}-p_{x_j}\partial_{\xi_j})\),
\[
 H_p=x_1\partial_{x_1}-\xi_1\partial_{\xi_1},
 \qquad H_p(\rho)=-R(\rho).
\]
In particular, the Hamilton field is a nonzero radial vector.

Let \(\kappa\) be a homogeneous symplectic local diffeomorphism. Since it
commutes with positive fibre dilations, differentiation of that
commutation relation gives \(\kappa_*R=R\). Symplectic covariance gives
\[
 \kappa_*H_p=H_{p\circ\kappa^{-1}}.
\]
Thus the transformed Hamilton field is still radial and nonzero at
\(\kappa(\rho)\). If the transformed defining symbol is further
multiplied by a smooth nonvanishing function \(a\), then, on its
characteristic hypersurface,
\[
 H_{a(p\circ\kappa^{-1})}
   =a\,H_{p\circ\kappa^{-1}},
\]
so radiality persists.

In contrast, for the coordinate symbol \(q(y,\eta)=\eta_2\),
\(H_q=\partial_{y_2}\). This is horizontal and nonzero, hence cannot
be proportional to the nonzero vertical radial vector at any
point of \(T^*\mathbb R^2\setminus0\), including points where
\(\eta_2=0\). The proposed homogeneous straightening is impossible.
\end{proof}

\begin{remark}
The two-dimensional example is the stabilization of the one-dimensional
symbol \(x\xi\). It avoids the incidental fact that the one-dimensional
coordinate characteristic set \(\eta=0\) misses the punctured cotangent
bundle. The obstruction uses three separately checked facts:
\(dp(\rho)\ne0\), homogeneity preserves the radial field, and a coordinate
symbol has a nonradial Hamilton field. It audits the canonical
transformation required by the usual elliptic FIO conjugation theorem;
it does not purport to rule out arbitrary unspecified kernel
factorizations. The missing nonradial condition is explicitly present
in the cited proposition of Duistermaat--Hörmander.
\end{remark}

\paragraph{Fourier sign used throughout.}
The reconstruction uses
\(\widehat u(\xi)=\int e^{-ix\xi}u(x)\,dx\).
Consequently \((x+i0)^{-1}\), the upper-half-plane boundary value of
\(1/z\), has positive, rather than negative, wavefront covectors.
The explicit positive-frequency representation and wavefront
calculation in
\citet[Appendix A, equation (A.4)]{bohrlefeuvrepaternain2024invariant}
independently checks this convention. P000584 does not specify its
Fourier sign, so its contrary excluded-half-ray wording is logged as
an ambiguity, not as a convention-independent mathematical error.
